\section{Developer experience with \enumo}
\label{sec:devexp}

In this section, we report on the developer experience of using
  \enumo to generate rulesets for the trigonometric domain used in \autoref{subsubsec:numbers}
  and the Halide domain used in \autoref{subsec:halideeval}.

\subsection{Developing custom workloads for trigonometry}
\label{subsec:trigexperience}

We describe the process of developing, in \enumo, the workloads
  we used in \autoref{subsubsec:numbers}.
The goal is to synthesize a set of rewrite rules that will perform well in
  \herbie, a rewrite-based tool for improving the accuracy of floating-point expressions.
We take an incremental approach, composing together the results of many
  invocations of the \ff algorithm.
At each step, rules learned in prior steps are used to shrink the \egraph
  during candidate generation and rule minimization.

\subsubsection{Constants}
First we enumerate terms that apply
 trigonometric functions to a set of constants.
\autoref{subsec:ffoverview} briefly describes this workload.
We explain the process in more detail here.
Specifically, we are interested in (possibly fractional) multiples of $\pi$.
The goal is to find identities like
  \C{(cos 0) \rewritesboth\, (sin (/ PI 2))}.
We are careful to filter out terms that are undefined like
  \C{(tan (/ PI 2))}: this term is equivalent to \C{(/ 1 0)} since
  \C{(sin (/ PI 2))} is 1 and \C{(cos (/ PI 2))} is 0.
Having undefined terms in the \egraph can cause other rules
  over rationals to discover unsound equivalences\footnote{For example,
  we include the rule \C{(/ x x) \rewritesto\, 1}, which is not sound if
  \C{x} is 0, but it is very useful when \C{x} is nonzero. To avoid
  unsoundness from this rule, it is important to prevent undefined
  terms from being represented in the \egraph at any point.}.

The benefit of learning these first is that they can help with additional
  constant folding.
Since these terms have no free variables
  but represent non-trivial identities,
  it is useful to discover them first in a separate phase.

The concrete workload contains 22 terms and is described by this \enumo program:
\begin{lstlisting}[name=trig]
  app = { (OP VAL) }
  ops = { sin cos tan }
  consts = { 0 (/ PI 6) (/ PI 4) (/ PI 3) (/ PI 2) PI (* PI 2) }
  W1 = app
          .plug("OP", ops)
          .plug("VAL", consts)
          .filter(Filter.Not(Filter.Contains("(tan (/ PI 2))")))
          .union({0 1})
\end{lstlisting}

The set of rewrite rules we find using \ff from this workload is:
\begin{lstlisting}[escapechar=@]
  (cos (/ PI 4)) @\rewritesboth@ (sin (/ PI 4))
  1 @\rewritesboth@ (tan (/ PI 4))
  1 @\rewritesboth@ (sin (/ PI 2))
  1 @\rewritesboth@ (cos (* PI 2))
  1 @\rewritesboth@ (cos 0)
  (tan PI) @\rewritesboth@ (cos (/ PI 2))
  (sin PI) @\rewritesboth@ (sin (* PI 2))
  (tan PI) @\rewritesboth@ (tan (* PI 2))
  (sin 0) @\rewritesboth@ (sin PI)
  (tan 0) @\rewritesboth@ (tan PI)
  0 @\rewritesboth@ (tan PI)
  (sin PI) @\rewritesboth@ (tan PI)
\end{lstlisting}

\subsubsection{Even/odd symmetry and periodicity}
Next, we create a workload to
  enumerate trigonometric functions whose arguments are transformed by (1) negation,
  (2) translation by $\pm \pi$, or (3) doubling.
The goal is to learn axioms like
  evenness: \C{f (- x) \rewritesboth\, f (x)},
  oddness: \C{f (- x) \rewritesboth\, -f (x)},
  and anti-periodicity: \C{f (+ x  P) \rewritesboth\, -f (x)}.
Evenness and oddness are useful for techniques like range reduction
  when constructing math libraries,
  where an even (or odd) function only needs
  to be implemented for half the range and
  the input (or output) can be transformed accordingly.
Conveniently, anti-periodicity axioms lead to periodicity identities:
  \C{f (+ x (* 2  P)) \rewritesto\, f (x)}.

Learning these axioms early is ideal.
Not only are they important properties of trigonometric functions,
  they serve as simplifying rewrites in later workloads,
  eliminating unnecessary negations and translations.

The concrete workload contains 50 terms and is described by this \enumo program:
\begin{lstlisting}
  simple_terms = app
                    .plug("OP", ops)
                    .plug("VAL", {a  (- a)  (+ PI a)  (- PI a)  (+ a a)})
  neg_terms = {(- X)}.plug("X", simple_terms)
  W2 = W1.union(simple_terms).union(neg_terms)
\end{lstlisting}

The set of rewrite rules we obtain from this workload is:
\begin{lstlisting}[escapechar=@]
  (- (cos a)) @\rewritesboth@ (cos (+ PI a))
  (sin a) @\rewritesboth@ (sin (- PI a))
  (tan a) @\rewritesboth@ (tan (+ PI a))
  (- (tan a)) @\rewritesboth@ (tan (- a))
  (- (sin a)) @\rewritesboth@ (sin (- a))
  (cos (- a)) @\rewritesboth@ (cos a)
\end{lstlisting}

\subsubsection{Sum-of-squares}
With this workload, we introduce sum-of-squares
  terms which have the form: \\
   \C{(+ (* (f x) (f x)) (* (g y) (g y)))}.

A number of trigonometric identities
  contain the square of a trigonometric function.
At this point, the full set of terms is too large,
so we prune away terms of the form \C{(f (+ x x))}.
With this workload, we learn the Pythagorean identity.

The concrete workload contains 116 terms and is described by this \enumo program:

\begin{lstlisting}
  no_trig_2x_filter = Filter.Not(Filter.Or(
                            Filter.Contains("(sin (+ x x))")
                            Filter.Contains("(cos (+ x x))")
                            Filter.Contains("(tan (+ x x))")))
  squares = { (sqr X) }.plug("X", app).plug("OP", ops).plug("VAL", {a b})
  add = {(+ E E) (- E E)}
  sum_of_squares = add.plug("E", squares)
  W3 = W2.filter(no_trig_2x_filter).union(sum_of_squares)
\end{lstlisting}

The set of rewrite rules \enumo finds from this workload is:
\begin{lstlisting}[escapechar=@]
  (+ (* (sin a) (sin a)) (* (cos a) (cos a))) @\rewritesto@ 1
  (- (* (cos b) (cos b)) (* (cos a) (cos a))) @\rewritesboth@ (- (* (sin a) (sin a)) (* (sin b) (sin b)))
  (- (* (sin b) (sin b)) (* (cos a) (cos a))) @\rewritesboth@ (- (* (sin a) (sin a)) (* (cos b) (cos b)))
\end{lstlisting}

\subsubsection{Coangles}
With this workload, we try to prove coangle identities.
Similar to the (anti-)periodicity identities,
  these identities allow us to eliminate translations
  by a quarter period.
Unlike in previous workloads where we simply added terms
  to the previous workload,
  here we start with a fresh set of more targeted terms
  to learn the identities quickly.
The concrete workload contains 13 terms and is described by this \enumo program:
\begin{lstlisting}
  nums = { -2 -1 0 1 2 }
  sin_cos = { sin cos }
  base = { a  (+ (/ PI 2) a)  (- (/ PI 2) a)  (* 2 a) }
  simple = app.plug("OP", sin_cos).plug("VAL", base)
  W4 = simple.union(nums)
\end{lstlisting}
The rewrite rule \enumo finds from this workload is:
\begin{lstlisting}[escapechar=@]
  (cos (- (/ PI 2) a)) @\rewritesboth@ (sin a)
\end{lstlisting}

\subsubsection{Power reduction}
With this workload,
  we try to learn power reduction identities,
  which rewrite trigonometric functions raised to a
  power into terms with smaller powers.
For example, the following identity is a power reduction identity:
  \C{(* (sin x) (sin x)) \rewritesto\, (/ (- 1 (cos (* 2 x))) 2)}.

The concrete workload contains 57 terms and is described by this \enumo program:
\begin{lstlisting}
shift = { X (- 1 X) (+ 1 X) }
scale_down = { X (/ X 2) }
shifted_simple = shift.plug("X", simple)
trivial_squares = { (sqr X) }.plug("X", app).plug("OP", sin_cos).plug("VAL", { a })
shifted_simple_sqrs = shifted_simple.union(trivial_squares)
scaled_shifted_sqrs = scale_down.plug("X", shifted_simple_sqrs)
W5 = scaled_shifted_sqrs.union(nums)
\end{lstlisting}

The set of rewrite rules \enumo finds from this workload is:
\begin{lstlisting}[escapechar=@]
  (/ (- 1 (cos (* 2 a))) 2) @\rewritesboth@ (* (sin a) (sin a))
  (/ (+ 1 (cos (* 2 a))) 2) @\rewritesboth@ (* (cos a) (cos a))
  (- a a) @\rewritesboth@ (- (- a a))  # minimization failed to eliminate this identity
\end{lstlisting}
Note that the last rule shows a limitation of \enumo's minimization approach: this
  identity should ideally have been eliminated, but due to the incomplete nature of minimization,
  it ends up in the final ruleset.

\subsubsection{Product-to-sum}
Similar to the previous workload,
  this workload explores identities with
  products of trigonometric functions such as
  \C{(* (sin x) (sin y))}.

To avoid enumerating terms from the previous workload,
 we filter out any squared terms.
These axioms are more general than the ones from the
  previous workload but are harder to learn, so we learn them separately.
While the more general versions together with other rewrites
  can in theory be used to derive the ones in the previous section,
  due to practical limitations like resource bounds,
  real-world equality saturation applications can benefit from both sets of
  rules.

The concrete workload contains 373 terms and is described by this \enumo program:

\begin{lstlisting}
non_square_filter = Filter.Not(Filter.Or(
          Filter.Contains("(* (sin x) (sin x))")
          Filter.Contains("(* (cos x) (cos x))")))
two_var = app
            .plug("OP", sin_cos)
            .plug("VAL", { a b (+ a b) (- a b) })
sum_two_vars = { (+ X Y) (- X Y) }
                  .plug("X", two_var)
                  .plug("Y", two_var)
prod_two_vars = { (* X Y) }
                   .plug("X", two_var)
                   .plug("Y", two_var)
                   .filter(non_square_filter)
sum_and_prod = sum_two_vars.union(prod_two_vars)
scaled_sum_prod = scale_down.plug("X", sum_and_prod)
W6 = scaled_sum_prod.union(nums)
\end{lstlisting}

The set of axioms \enumo learns from this workload is:
\begin{lstlisting}[escapechar=@]
  (* (sin b) (sin a)) @\rewritesboth@ (/ (- (cos (- b a)) (cos (+ b a))) 2)
  (* (cos b) (cos a)) @\rewritesboth@ (/ (+ (cos (+ b a)) (cos (- b a))) 2)
  (+ b a) @\rewritesto@ (+ a b) # minimization failed to eliminate this identity
\end{lstlisting}

\subsubsection{Sums}
Finally, we find identities involving trigonometric
  functions with sums as their arguments, such as
  \C{(sin (+ x y))}.

These are related to the product-to-sum workload
  since products of trigonometric functions rewrite
  into trigonometric functions of sums and vice versa.
This workload enumerates many terms, so we impose a number
of filters, including removing double-angle terms,
  e.g., \C{(sin (+ x x))}, and squared terms, e.g., \C{(sin (* x x))},
  to keep the search tractable.

The concrete workload contains 287 terms and is described by this \enumo program:

\begin{lstlisting}
trig_no_sub_filter = Filter.Not(Filter.Or(
            Filter.Contains("(cos (- a b))")
            Filter.Contains("(sin (- a b))")))
two_x_filter = Filter.Not(Filter.Contains("(+ x x)"))
two_var_no_sub = two_var.filter(trig_no_sub_filter)
sum_of_prod = { (+ X Y) (- X Y) }
              .plug("X", prod_two_vars)
              .plug("Y", prod_two_vars)
              .filter(two_x_filter)
              .filter(non_square_filter)
W7 = two_var_no_sub.union(sum_of_prod).union(nums)
\end{lstlisting}

The set of axioms \enumo learns from this workload is:
\begin{lstlisting}[escapechar=@]
  (- (* (cos b) (cos a)) (* (sin b) (sin a))) @\rewritesboth@ (cos (+ b a))
  (sin (+ b a)) @\rewritesboth@ (+ (* (sin b) (cos a)) (* (sin a) (cos b)))
\end{lstlisting}

Note that the incremental approach to rule inference is crucial.
Rules learned at each step are used as prior rules in subsequent steps.
It is possible to run \ff with a single workload composed of all of the
  terms from each of the 7 workloads described above, but the rules synthesized
  contain more duplicates and overly complicated rules
  (e.g., \C{(- (- b a) (- b a)) \rewritesto\, (sin PI)} instead of \C{0 \rewritesto\, (sin PI)}).
With \enumo, it is easy to design small workloads and compose
  rulesets incrementally.

\subsection{Developing custom workloads for a Halide-inspired domain}
\label{subsec:halide-appendix}
It is possible to \textit{overfit} an \enumo workload by only enumerating terms
  that correspond to the left- and right-hand sides of a desired rule.
For example, creating a workload consisting only of the terms
  \C{(- (* x y) (* z y))} and \C{(* (- x z) y)} in
  order to synthesize the rule \C{(- (* x y) (* z y)) \rewritesboth\, (* (- x z) y)}
  is possible, but requires that the desired rule is known beforehand.
However, the real value of \enumo is its ability to discover useful rewrite rules
  \textit{without} the domain expert needing to already know what rules they
  are looking for.
\enumo allows domain experts to guide the term enumeration, enabling \enumo
  to scale past tools that rely on exhaustive term enumeration.
While it is possible to use guided term enumeration to find specific desired rules,
  as described above, it is more useful to use \enumo workloads to encode
  general insights about the domain of interest.
To see how \enumo enables users to encode \textit{broad}, rather than specific,
  domain knowledge, consider the \enumo program for the Halide-inspired
  grammar (\autoref{subsec:halideeval}).
Here, we do not guide the enumeration towards specific rules we want to synthesize,
  but rather we encode intuition about which parts of the search space are worth
  exploring.

First, we group the operators in the domain into semantic categories:
  boolean operators, arithmetic operators, and predicate operators.
Then, we compute an \enumo workload corresponding to exhaustive enumeration
  using each subset of operators, and synthesize rules using each workload separately.
The intuition guiding this strategy is that useful rules are more likely to be found
  within each category of operators than between categories.
Of course, there are also useful rewrite rules that use operators from multiple categories,
  so we also create a workload that enumerates over all of the operators.
However, we are able to use all of the rules we learned previously to shrink the
  \egraph and make rule inference more tractable.
Finally, we enumerate workloads with larger terms, but apply the \enumo{} \textsf{Canon}
  filter to remove terms whose variables do not appear in canonical order,
  again making the search space smaller and more tractable.
Since we have already learned rules that enable reordering
  and re-associating terms, variable order within a term should not matter as much.

The domain expertise we utilize to develop this \enumo program is
  quite general, not overly specific to the rules we intend to synthesize.
However, it is still valuable to guide the search space in broad ways,
  which is not feasible with one-shot theory explorers like \ruler.
By separating the rule inference problem into smaller searches over subsets of
  the enumeration space, we are able to synthesize a more powerful ruleset
  than prior work.
